\documentclass[../main.tex]{subfiles}
\begin{document}

\section{Code Implementation Details}
To ensure reproducibility and clarity, this section details the PyTorch implementation of the proposed Cross-modal Circle Loss and its integration into the training loop with the Curriculum Hard-Mining schedule.

\subsection{Circle Loss Construction}
\noindent The core mathematical formulation of Circle Loss is encapsulated in the function \texttt{compute\_cross\_modal\_circle}. Unlike standard implementations that operate on a single modality, our implementation explicitly constructs a cross-modal similarity matrix between image embeddings and text embeddings. The implementation follows a vectorized approach for efficiency:

% Thêm tùy chọn itemsep và topsep để danh sách trông thoáng hơn
\begin{enumerate}[itemsep=0.5em, topsep=0.5em]
    \item \textbf{Normalization:} Both image and text features are L2-normalized to ensure they lie on the hypersphere.
    \item \textbf{Similarity Matrix:} I compute the cosine similarity matrix $S = I \cdot T^T$.
    \item \textbf{Masking:} I utilize the identity labels (PIDs) to create boolean masks for positive pairs ($s_p$) and negative pairs ($s_n$).
    \item \textbf{Adaptive Re-weighting:} The weighting factors $\alpha_p$ and $\alpha_n$ are computed dynamically based on the deviation of the similarity scores from the optimum boundaries ($1+m$ and $-m$, respectively).
    \item \textbf{Logit Scaling:} The re-weighted scores are scaled by the inverse temperature $\gamma$ (Gamma) before being passed to the \textbf{Softplus} activation function, which serves as a smooth approximation of the ReLU function in the final loss calculation.
\end{enumerate}

% --- CODE BLOCK 1 ---
% Sử dụng basicstyle=\ttfamily\scriptsize để giảm cỡ chữ
% breaklines=true và breakatwhitespace=false để ép xuống dòng các dòng dài
\begin{lstlisting}[
    language=Python, 
    caption={PyTorch implementation of Cross-modal Circle Loss},
    label={lst:circle_loss_impl},
    basicstyle=\ttfamily\scriptsize, % Chữ nhỏ hơn để tránh tràn
    breaklines=true,                 % Cho phép ngắt dòng
    breakatwhitespace=false,         % Ngắt cả ở những chỗ không phải khoảng trắng
    xleftmargin=1.5em,               % Lùi lề trái vào một chút
    xrightmargin=0.5em               % Lùi lề phải vào một chút
]
def compute_cross_modal_circle(image_features, text_features, pids, m=0.25, gamma=128):
    """
    Circle Loss between 2 modalities
    """
    image_features = F.normalize(image_features, dim=1, p=2)
    text_features = F.normalize(text_features, dim=1, p=2)

    sim_mat = torch.matmul(image_features, text_features.t())

    pids = pids.view(-1, 1)
    pos_mask = torch.eq(pids, pids.t()).float()
    neg_mask = 1 - pos_mask

    s_p = sim_mat[pos_mask.bool()]
    s_n = sim_mat[neg_mask.bool()]

    if s_p.numel() == 0 or s_n.numel() == 0:
        # Ensure gradient flow even if no positive/negative pairs exist in batch
        return torch.tensor(0.0, device=image_features.device, requires_grad=True)

    # Adaptive re-weighting
    alpha_p = torch.clamp_min(-s_p.detach() + 1 + m, min=0.)
    alpha_n = torch.clamp_min(s_n.detach() + m, min=0.)

    delta_p = 1 - m
    delta_n = m

    # Logit scaling
    logit_p = - gamma * alpha_p * (s_p - delta_p)
    logit_n = gamma * alpha_n * (s_n - delta_n)

    # Softplus activation as smooth approximation of ReLU
    soft_plus = torch.nn.Softplus()
    loss = soft_plus(
        torch.logsumexp(logit_p, dim=0) +
        torch.logsumexp(logit_n, dim=0)
    )
    return loss
\end{lstlisting}

\subsection{Circle Loss Implementation to Model}
The integration of this objective is handled within the \texttt{forward} method of the \textbf{TBPS} class (in \texttt{tbps.py}). This method orchestrates the hybrid optimization strategy, managing the flow of features through the backbone and applying the Curriculum Hard-Mining Schedule.

\noindent \textbf{Curriculum Weight Scheduling:} A critical component of our implementation is the dynamic adjustment of the Circle Loss weight ($\alpha_c$). Instead of a static hyperparameter, I implement a linear warmup logic directly inside the training loop.

\begin{itemize}[itemsep=0.5em, topsep=0.5em]
    \item \textbf{Warm-up Phase ($t < T_{warm up}$):} The weight is clamped to 0.0, effectively disabling Circle Loss to allow N-ITC to stabilize global alignment.
    \item \textbf{Ramp-up Phase:} The weight increases linearly from 0.0 to the target value ($\alpha_{target} = 0.1$) over a defined period ($T_{ramp-len}$), allowing the model to gradually adapt to the harder geometric constraints.
\end{itemize}

\noindent \textbf{Multi-View Supervision (MVS) Integration:} To enforce robustness, the Circle Loss is applied in a dual-branch manner:

\begin{enumerate}[itemsep=0.5em, topsep=0.5em]
    \item \textbf{Primary Branch:} Computes loss between the original image features and text features.
    \item \textbf{MVS Branch:} Computes loss between augmented image features (e.g., random erasing, cropping) and text features. The final \texttt{circle\_loss} is the average of these two branches, weighted by the current curriculum factor.
\end{enumerate}

% --- CODE BLOCK 2 ---
% Block này rất dài, cần scriptsize và setting ngắt dòng mạnh
\begin{lstlisting}[
    language=Python, 
    caption={Integration of Circle Loss and MVS in training loop},
    label={lst:training_loop_impl},
    basicstyle=\ttfamily\scriptsize, % Sử dụng scriptsize để vừa vặn hơn
    breaklines=true,
    breakatwhitespace=false,
    xleftmargin=1.5em, 
    xrightmargin=0.5em
]
# --- E2. AUXILIARY CROSS-MODAL (Combined Static + Curriculum) ---
if strategy in ["auxiliary_cross", "auxiliary_cross_curriculum"]:

    # 1. Calculate N-ITC Loss first (Global Alignment)
    if self.config.loss.get("NITC", None):
        sim_targets = self.prepare_sim_targets(batch["pids"], self.use_sigmoid)

        # Detach features if alpha is 0 to stop gradient backprop
        image_pooler_output_stopped = (
            image_pooler_output.clone().detach() if alpha != 0 else None
        )
        caption_pooler_output_stopped = (
            caption_pooler_output.clone().detach() if alpha != 0 else None
        )

        nitc_loss = objectives.compute_constrative(
            image_features=image_pooler_output,
            text_features=caption_pooler_output,
            image_features_stopped=image_pooler_output_stopped,
            text_features_stopped=caption_pooler_output_stopped,
            sim_targets=sim_targets,
            alpha=alpha,
            logit_scale=logit_scale,
            logit_bias=self.backbone.logit_bias,
            use_sigmoid=self.use_sigmoid,
            weights=weights,
        )

        # Apply MVS to N-ITC if enabled
        if self.config.loss.get("MVS", None):
            aug_images = batch["aug_images"]
            aug_images_features = self.encode_image(aug_images)
            aug_images_features_stopped = (
                aug_images_features.clone().detach() if alpha != 0 else None
            )
            # Compute contrastive loss for augmented view
            augmented_nitc_loss = objectives.compute_constrative(
                image_features=image_pooler_output, 
                text_features=caption_pooler_output,
                image_features_stopped=aug_images_features_stopped,
                text_features_stopped=caption_pooler_output_stopped,
                sim_targets=sim_targets,
                # ... (other parameters omitted for brevity) ...
                weights=weights,
            )
            # Average primary and augmented losses
            nitc_loss = (nitc_loss + augmented_nitc_loss) / 2
        
        ret.update({"nitc_loss": nitc_loss * self.config.loss.nitc_loss_weight})

    # 2. Calculate Circle Loss (Geometric Constraints) based on schedule
    if self.config.loss.get("CIR", None) and current_circle_weight > 0:
        circle_m = self.config.loss.get("circle_margin", 0.25)
        circle_gamma = self.config.loss.get("circle_gamma", 128)

        # Primary Branch: Cross-Modal Circle Loss
        cm_circle_loss = objectives.compute_cross_modal_circle(
            image_features=image_pooler_output,
            text_features=caption_pooler_output,
            pids=batch["pids"],
            m=circle_m,
            gamma=circle_gamma
        )

        final_circle_loss = cm_circle_loss
        
        # MVS Branch for Circle Loss
        if self.config.loss.get("MVS", None):
            # Ensure augmented features are encoded if not already done
            if 'aug_images_features' not in locals():
                aug_images = batch["aug_images"]
                aug_images_features = self.encode_image(aug_images)

            aug_cm_circle_loss = objectives.compute_cross_modal_circle(
                image_features=aug_images_features,
                text_features=caption_pooler_output,
                pids=batch["pids"],
                m=circle_m,
                gamma=circle_gamma
            )
            # Average the losses
            final_circle_loss = (cm_circle_loss + aug_cm_circle_loss) / 2

        # Apply curriculum weight
        ret.update({"circle_loss": final_circle_loss * current_circle_weight})

    elif self.config.loss.get("CIR", None):
        # If Circle Loss is enabled but weight is 0 (warm-up phase), return 0 loss with grad
        ret.update({"circle_loss": torch.tensor(0.0, device=image_pooler_output.device, requires_grad=True)})
\end{lstlisting}

\end{document}